\scanpage{78}
Ainsi on a prouvé

\uline{Prop} Supposons $2$ régulier dans $A$, $A_0=A/2A$ réduit et \unsure{intégralement} clos, $\Delta\in A$ tel que $\Delta_0\in A_0$ soit él. régulier de $A_0$. Soit $\unsure{k=2d-|*=2d-|k=2d-1|\nu=2d+1}$, comme nous notons supra. Pour qu'il existe $b',c'\in A$, \unsure{ou on|ou un} \inserted{$\unsure{v|\upsilon|u}\in A^{\deleted{\unsure{*}}}$} a
% slanted
\margin{\uline{NB} -- \unsure{Lieu|au lieu} de $v\in A^*$, il suffit existe $v\in A$ tel que $v$ soit régulier ds $A$ et $v_0$ régulier\unsure{!?|?} ds $A_0$}
\[
\left\{\begin{aligned}
{b'}^2-4c'&=\Delta^{\deleted{\unsure{*}}}v^2\\
v_0&\text{ est $A_0$-régulier}
\end{aligned}\right.
\]
\unsure{faut s.|f. et s.} qu'il existe $b,c\in A$ tel que
\[
 b^2-4c=\Delta.
\]

\deleted{Supposons}\inserted{Considérons} maintenant \deleted{que \unsure{l'on ait|l'anneau}}\inserted{l'anneau}
\[
 A_{2\Delta}\deleted{\unsure{=A}}\inserted{=(A_2)_\Delta=(A_\Delta)_2},\ \text{localisé de $A$ pour $2\Delta$} \tag*{(\unsure{a|$\alpha$})}
\]
et l'Algèbre
\[
 B_{2\Delta}=A_{2\Delta}[T]/(T^2-\Delta)
\]
obtenue en extrayant \deleted{la} racine carrée de $\Delta$. \deleted{Suppos}\inserted{C'est une $A_{2\Delta}$-alg. étale galoisienne de grpe $\mathbb Z/2$.} \uline{Supposons} que $B_{2\Delta}$ se prolonge en une Algèbre $B_\Delta$ sur $A_\Delta$ \inserted{finie} étale \deleted{\unsure{de}} [qui sera unique : \unsure{isom|sens} unique près, car $2$ est régulier dans $A_\Delta\ldots$]. Supposons de plus \unsure{que|q.|Z}
\[
 \text{\textcircled{e}}\quad \boxed{\operatorname{Pic} A_\Delta=0}
\]
alors on sait (comme $B_\Delta$ est finie loc libre de rang $2$) que l'on peut écrire

% archive sheet number: 77
